\documentclass[11pt,a4paper]{report}

%INPUT of SETTINGS (usepackage, styles, new command,...)
\usepackage{geometry}
\geometry{tmargin=2cm,bmargin=2cm,lmargin=3cm,rmargin=3cm,includefoot, includehead}
\usepackage[english]{babel}
\usepackage[T1]{fontenc}
\usepackage[latin1]{inputenc}
\usepackage{graphicx}
\usepackage{titlesec, blindtext, color}
\usepackage{fancyhdr}
\usepackage{hyperref}
\usepackage{float}
\usepackage{enumitem}
\usepackage{lastpage}
\usepackage{listings}
\usepackage{color}
\usepackage{amsmath}
\usepackage{tabularx}
\usepackage{array}
\usepackage{verbatim}
\usepackage{booktabs}
\usepackage{ltxtable} 
\usepackage{multirow}
\makeatletter\let\MPtrue\@minipagetrue\makeatother

%----------------------------------------------------------------------------------------
%   Title Style Chapter
%----------------------------------------------------------------------------------------
\definecolor{gray75}{gray}{0.70}
\newcommand{\hsp}{\hspace{15pt}}
\newcommand\tab[1][1cm]{\hspace*{#1}}
\newcommand{\forceindent}{\leavevmode{\parindent=1em\indent}}
\titleformat{\chapter}[hang]{\Huge\bfseries}{\thechapter\hsp\textcolor{gray75}{|}\hsp}{0pt}{\Huge\bfseries}


%----------------------------------------------------------------------------------------
%  ToC
%----------------------------------------------------------------------------------------
\hypersetup{
    colorlinks,
    citecolor=black,
    filecolor=black,
    linkcolor=black,
    urlcolor=black
}

% Renaming of ToC
\addto\captionsenglish{% Replace "english" with the language you use
  \renewcommand{\contentsname}%
    {Table of Contents}%
}

%----------------------------------------------------------------------------------------
%   Header/Footer
%----------------------------------------------------------------------------------------
\pagestyle{fancy} 
\fancyhf{}
%\renewcommand{\chaptermark}[1]{\markboth{#1}{}}

%\renewcommand{\chaptermark}[1]{%
% \markboth{\chaptername
% \ \thechapter.\ #1}{}}

%\renewcommand{\sectionmark}[1]{\markright{\thesection.\ #1}}

\lhead{\nouppercase{\leftmark}}
%\chead{}
\rhead{\nouppercase{\rightmark}}
\lfoot{DD - \textit{Version 1.0}}
%\cfoot{}
\rfoot{\thepage}
\renewcommand{\headrulewidth}{0.008pt}
\renewcommand{\footrulewidth}{0.008pt}

% Redefine the plain page style
\fancypagestyle{plain}{%
	\fancyhf{}%
	\fancyfoot[R]{\thepage}%
	\renewcommand{\headrulewidth}{0pt}% Line at the header invisible
	\renewcommand{\footrulewidth}{0.008pt}% Line at the footer visible
}

%-----------
% New Page for each section
%---------
%\newcommand{\sectionbreak}{\clearpage}

%-------
% Numbered list for goals
%-------
\newlist{Goal}{enumerate}{1}
\setlist[Goal]{label=[G\arabic*],leftmargin=0.6in}

%-------
% Numbered list for requirements
%-------
\newlist{Req}{enumerate}{1}
\setlist[Req]{label=[R\arabic*],leftmargin=0.6in}

%-------
% Numbered list for domain assumptions
%-------
\newlist{DA}{enumerate}{1}
\setlist[DA]{label=[DA\arabic*],leftmargin=0.6in}

%-------
% Numbered list for text assumptions
%-------
\newlist{TA}{enumerate}{1}
\setlist[TA]{label=[TA\arabic*],leftmargin=0.6in}

%-----
% Image box, arguments: file path, caption
%----
\newcommand{\imagebox}[3]{
\begin{figure}[H]
\centering
\fbox{\includegraphics[width=\textwidth,keepaspectratio,angle=#3]{#1}}
\caption{#2}
\end{figure}
}

\newcommand{\imageboxNoFit}[4]{
\begin{figure}[H]
\centering
\fbox{\includegraphics[keepaspectratio,angle=#3,scale=#4]{#1}}
\caption{#2}
\end{figure}
}

%---
% Alloy listing
%---
\definecolor{pgreen}{rgb}{0,0.64,0.013}
\definecolor{pblue}{rgb}{0.13,0.13,1}

\lstdefinelanguage{pseudocode}{
	keywords={
		function, if, endif, for, each, endfor, return, called, in, boolean, integer, float, Position, Car, List, PositionMediaVariance, not, break},
	literate=%
	{:}{$\colon$}1
	%{|}{$\bullet$}1
	%{==}{$=$}1
	%{=}{$=$}1
	{!=}{$\neq$}1
	%{&&}{$\land$}1
	%{||}{$\lor$}1
	{<=}{$\le$}1
	{>=}{$\ge$}1
	%{all}{$\forall$}1
	%{exists}{$\exists$}1
	{!in}{$\not\in$}1
	{\\in}{$\in$}1
	{=>}{$\implies$}3
	% the following isn't actually Alloy, but it gives the option to produce nicer latex
	{|=>}{$\Rightarrow$}2
	{<=set}{$\subseteq$}1
	{+set}{$\cup$}1
	{*set}{$\cap$}1
	{==>}{$\Longrightarrow$}3
	{<=>}{$\iff$}4
	{...}{$\ldots$}1
	{\\hl}{$\hline$}1
	{\\alpha}{$\alpha$}1
	{\\beta}{$\beta$}1
	{\\gamma}{$\gamma$}1
	{\\delta}{$\delta$}1
	{\\epsilon}{$\epsilon$}1
	{\\zeta}{$\zeta$}1
	{\\eta}{$\eta$}1
	{\\theta}{$\theta$}1
	{\\iota}{$\iota$}1
	{\\kappa}{$\kappa$}1
	{\\lambda}{$\lambda$}1
	{\\mu}{$\mu$}1
	{\\nu}{$\nu$}1
	{\\xi}{$\xi$}1
	{\\pi}{$\pi$}1
	{\\rho}{$\rho$}1
	{\\sigma}{$\sigma$}1
	{\\tau}{$\tau$}1
	{\\upsilon}{$\upsilon$}1
	{\\phi}{$\phi$}1
	{\\chi}{$\chi$}1
	{\\psi}{$\psi$}1
	{\\omega}{$\omega$}1
	{\\Gamma}{$\Gamma$}1
	{\\Delta}{$\Delta$}1
	{\\Theta}{$\Theta$}1
	{\\Lambda}{$\Lambda$}1
	{\\Xi}{$\Xi$}1
	{\\Pi}{$\Pi$}1
	{\\Sigma}{$\Sigma$}1
	{\\Upsilon}{$\Upsilon$}1
	{\\Phi}{$\Phi$}1
	{\\Psi}{$\Psi$}1
	{\\Omega}{$\Omega$}1
	{\\EOF}{\;}1,
	sensitive=true,
	morecomment=[l]//,%
	%morecomment=[l]{--},%
	morecomment=[s]{/*}{*/},%
	morestring=[b]",
	numbers=none,
	firstnumber=1,
	%numberstyle=\tiny,
	%basicstyle=\scriptsize\ttfamily,
	commentstyle=\color{pgreen},
	keywordstyle=\color{pblue},
	%ndkeywordstyle=\bfseries,
}

\lstnewenvironment{pseudocodeEnv}{%	
	\lstset{language=pseudocode,
		floatplacement={tbp},
		captionpos=b,
		xleftmargin=2pt,
		xrightmargin=2pt,
		breaklines=true,
		tabsize=3,
		%basicstyle=\ttfamily,
	}
} {}
%--- End of listing for Alloy
%
%---
% Recall item commands
%---
\newcommand\itemtext[2]{%
	\expandafter\gdef\csname item#1\endcsname{#2}%
	\label{#1}#2}

\newcommand\useitemtext[2]{\csname item#1\endcsname}

%--- 
% Table for use case descriptions
%---
\newcommand{\usecasedesc}[6]{%
	\begin{table}[H]
		\begin{tabular}{p{0.3\linewidth} p{0.7\linewidth}}
			%\begin{tabular}{\linewidth}{>{\setlength\hsize{.25\hsize}}X>{\setlength\hsize{.75\hsize}\setlength\linewidth{\hsize}}X}
			\textit{Name} & #1 \\
			\textit{Actors} & #2 \\
			\textit{Entry conditions} &\MPtrue#3\\
			\textit{Flow of events} &\MPtrue#4\\
			\textit{Exit conditions} & #5\\
			\textit{Exceptions} & #6\\
		\end{tabular}
	\end{table}
}

%\newcommand{optusecase}{[leftmargin=0.50cm,nosep]}


%\setlength{\abovecaptionskip}{3pt}

%-------------------------------------------------------------------------
%   START DOCUMENT 
%-------------------------------------------------------------------------
\begin{document}

% INPUT OF TITLE PAGE
%----------------------------------------------------------------------------------------
%	TITLEPAGE
%----------------------------------------------------------------------------------------
\begin{titlepage}

\newcommand{\HRule}{\rule{\linewidth}{0.3mm}} % Defines a new command for the horizontal lines, change thickness here

%\vspace*{0.3cm}
 
\center % Center everything on the page

%----------------------------------------------------------------------------------------
%	LOGO SECTION
%----------------------------------------------------------------------------------------
\includegraphics[scale=0.33]{../../logo.png}\\[0.4cm] 

%----------------------------------------------------------------------------------------
%	HEADING SECTION
%----------------------------------------------------------------------------------------
\textsc{\Large \textbf{Politecnico di Milano}\\MSc Computer Science and Engineering}\\[1.7cm] \textsc{\Large \textbf{Software Engineering 2}\\[0.1cm]Academic Year 2016-2017}\\[1cm] 
%\textsc{\LARGE Minor Heading}\\[0.5cm] 

%----------------------------------------------------------------------------------------
%	TITLE SECTION
%----------------------------------------------------------------------------------------

\HRule \\[0.4cm]
{ \huge \bfseries Integration Test Plan Document\\[0.5cm]\textit{PowerEnJoy}}\\[0.4cm] % Title of document
\HRule \\[1.5cm]
 
%----------------------------------------------------------------------------------------
%	AUTHOR SECTION
%----------------------------------------------------------------------------------------
\Large \emph{Authors:}\\
\LARGE{
\begin{center}
\begin{tabular}{ r  l }
	Melloni Giulio & 876279\\
	Renzi Marco & 878269\\
	Testa Filippo & 875456 \\
\end{tabular}
\end{center}
}
\vspace*{0.8cm}
\Large{\emph{Reference Professor:} \\
{\textsc{Mottola} Luca} 
}

%----------------------------------------------------------------------------------------
%	DATE SECTION
%----------------------------------------------------------------------------------------
 \vspace*{1cm}
{\LARGE {\it Release Date: January 15$^{th}$, 2017\\Version 1.0}}
%vfill % Fill the rest of the page with whitespace

\end{titlepage}

%-------------------------------------------------------------------------
%   TOC
%-------------------------------------------------------------------------
%\pagestyle{plain}
\pagenumbering{roman}
\thispagestyle{empty}
\tableofcontents
\listoffigures
\cleardoublepage
\pagenumbering{arabic}
\pagestyle{fancy}

%--------------------------------------------------------------------
%  BEGIN TEXT
%--------------------------------------------------------------------
\chapter{Introduction}
\section{Revision History}
%Record all revisions to the document
\begin{tabularx}{\textwidth}{|l|l|X|}
	\hline
	Version & Date & Summary\\
	\hline
	1.0 & 15/01/2017 & First release of the document\\
	\hline
\end{tabularx}

\section{Purpose and Scope}
The Integration Test Plan Document (after referred simply as \textit{ITPD}) covers the whole process of integration of the components of the system; in particular it explains in a detailed way how these ones will be integrated and the specific sequence that must be followed.\\
The document firstly tackles the integration plan issue starting from some assumptions, as listed in the Entry Criteria section; after that the component diagram presented in the DD is splitted up in different macroareas that will be useful to better realize the bottom-up approach (Elements to be integrated section).\\
Then the strategy that leads the integration plan is stated and, obviously along with it, its rationale. Following the strategy, the sequence of integration of the components and the single steps are derived.\\
The main goal of the integration test plan is to ensure the all the developed components properly combine together in order to realize all the functionalities and also to check that any unexpected behaviours of the system arise.
\section{List of Definitions and Abbreviations}
\begin{description}
	\item[PowerEnJoy]is the name of the system that has to be developed.
	\item[System]sometimes called also \textit{system-to-be}, represents the application that will be described and implemented.
	In particular, its structure and implementation will be explained in the following documents. People that will use the car-sharing service will interact with it, via some interfaces, in order to complete some operations (e.g.: reservation and renting).
	\item[Renting]it is the act of picking-up an available car and of starting to drive.
	\item[Ride] the event of picking-up a car, driving through the city and parking it. Every Ride is associated to a single user and to a single car.
	\item[Reservation]it is the action of booking an available car.
	\item[Car] a car is an electrical vehicle that will be used by a registered user.
	\item[Not Registered User] indicates a person who hasn't registered to the system yet; for this reason he can't access to any of the offered function. The only possible action that he can carry out is the registration to get a personal account.
	\item[Registered User] interacts with the system to use the sharing service. He has an account (which contains personal information, driving license number and payment data) that must be used to access to the application in order to exploit all the functionalities.
	\item[Employee] it's a person who works for the company, whose main task is to plug into the power grid those cars that haven't been plugged in by the users. He is also in charge of taking care of the status of the cars and of moving the vehicles from a safe area to a charging area and vice versa if needed.
	\item[Safe Area] indicates a set of parking lots where the users have to leave the car at the end of the rent; the set of the Safe Areas is pre-defined by the system management. These areas are spread all over the city.
	\item[Plug] defines the electrical component that physically connects the car to the power grid.
	\item[Charging Area] is a special \textit{Safe Area} that also provides a certain number of plugs that connect the cars to the power grid in order to recharge the battery.
	\item[Registration] the procedure that an unregistered user has to perform to become a registered user. At the end, the unregistered user will have an account. To complete this operation three different types of data are required: personal information, driving license number and payment info.
	\item[Search] this functionality lets the registered user search for available cars within a certain range from his/her current position or from a specified address.
	\item[RASD] is the acronym of \textit{Requirements Analysis and Specification Document}
	\item[DD] is the acronym of \textit{Design Document}
	\item[ITPD] is the acronym of \textit{Integration Test Plan Document}
\end{description}
\section{List of Reference Documents}
	\begin{itemize}
		\item Project Assignments 2016-2017
		\item RASD v1.1
		\item DD v1.0
		%TODO documentation of used tools
	\end{itemize}
\chapter{Integration Strategy}
\section{Entry Criteria}
There are some criteria that impose some conditions on the project testing phase.\\
Firstly some considerations on the level of completion of the components with respect to their functionalities:
\begin{itemize}
	\item The \textbf{Dispatcher} must have been fully implemented in order to route the simulated requests
	\item Controllers like the \textbf{ReservationManager}, the \textbf{RegistrationManager}, the \textbf{StateManager}, the \textbf{LogInManager} and the \textbf{RideManager} have to expose sufficiently developed interfaces in order to be able to test the requests management
	\item No specific constraints on the \textbf{ViewRender} for the first structural testing phase which involves operations that don't integrate the UI
	\item Components like the \textbf{Payment Manager} and the \textbf{MapController} that are to be linked with third-part components (\textbf{Payment System} and \textbf{MapServce}) must have fully developed in order to use the external APIs
\end{itemize}
Secondly, the \textit{Requirements Analysis and Specification Document} and the \textit{Design Document} must have been written.\\
Thirdly, the components must have been individually tested (unit testing is not part of this testing phase) in order to ensure that bugs from the upcoming integration tests will be caused exclusively by the iterations among these components and not by any kind of internal problem.
%Specify the criteria that must be met before integration testing of specific elements may begin (e.g., functions must have been unit tested).
\section{Elements to be Integrated}
%Identify the components to be integrated, refer to your design document to identify
% such components in a way that is consistent with your design.
In order to build the full \textit{PowerEnJoy} system all its components have to be properly integrated. In this section the focus is on which components are selected and how these are aggregated.\\Let us consider the component diagram of the \textit{Design Document} to refer to the components to be integrated. For the integration testing purpose it is useful to organize the components into logical \textbf{Macro Areas} that will support the testing process as explained in the \textit{Integration Test Strategy} section:
\begin{itemize}
\item \textbf{Input Area} includes \textit{ViewRender} and \textit{Dispatcher} components. This pair of modules should be tested together to ensure that all input requests are properly received by the system.
\imageboxNoFit{../img/input}{Input Area}{0}{0.3}
\item \textbf{Management Area} includes the \textit{ReservationManager}, the \textit{RegistrationManager}, the \textit{StateManager}, the \textit{LoginManager}, the \textit{MapController}, the \textit{RideManager} and the \textit{Dispatcher}. These modules are responsible for the business logic of the application and consequently should be tested together.
\imageboxNoFit{../img/management}{Management Area}{0}{0.3}
\item \textbf{Render Area} is the set made of the \textit{ViewRender} and of the \textit{ReservationManager}, the \textit{RegistrationManager}, the \textit{StateManager}, the \textit{LoginManager}, the \textit{MapController}, the \textit{RideManager}. This logical area has to be tested in order to ensure that all managers can update the view of the application without bugs.
\imageboxNoFit{../img/render}{Render Area}{0}{0.3}
\item \textbf{Ride Area} includes \textit{RideManager},  \textit{MapController}, \textit{RideCostCalculator} and \textit{PaymentManager}. The tests on this area is crucial because it is responsible of the costs computation and of the payment process.
\imageboxNoFit{../img/ride}{Ride Area}{0}{0.3}
\item \textbf{Data Area} is the group of components that deal with the \textit{Model} and the external \textit{DBMS}. This is made of the \textit{Model} itself and of the \textit{ReservationManager}, the \textit{RegistrationManager}, the \textit{StateManager}, the \textit{LoginManager}, the \textit{MapController}, the \textit{RideManager}. Tests in this area aims at verifying the correctness of data through the various operations that the system has to perform on them.
\imageboxNoFit{../img/data}{Data Area}{0}{0.3}
\item \textbf{CarCommunication Area} is the pair of \textit{ServerCommunicationManager}and \textit{CarCommunicationManager}. Here the tests have to ensure that flow of information in both directions is feasible and consistent.
\imageboxNoFit{../img/carCommunication}{CarCommunication Area}{0}{0.3}
\item \textbf{Car Area} is the logical set of components that have to be tested on the car. \textit{CarCommunicationManager}, \textit{CentralUnit} and \textit{ScreenManager} are part of the Built-in sw for the car. 
\imageboxNoFit{../img/car}{Car Area}{0}{0.3}
\end{itemize}
Please note that the given groupings do not represent a partition of the set of components of the system (some components are shared by more than one macro area) but just a logical division that is convenient to carry out the integration testing. 
Finally, a remark on the external components (\textit{MapService}, \textit{Payment System} and \textit{DBMS}): they are already available for integration testing and they only require a sufficient level of completion of the internal components to be actually tested.
\section{Integration Test Strategy}
%Describe the integration testing approach (top-down, bottom-up, functional groupings,
%etc.) and the rationale for the choosing that approach.
The integration testing process will be carried out with both bottom-up and threads approaches. In particular, the bottom-up testing will be executed between modules that belong to the same macro area (as defined in the \textit{Elements to be Integrated} section) throughout their development process while the threads analysis will be eventually performed among modules of different areas when the previous internal tests are successfully passed.\\
This testing strategy is incremental by construction because it follows the the development of the components and consequently it makes it easier to spot possible errors during the implementation. As portions of components are added to the existing ones, the integration testing will be triggered on the new parts making use of suitable drivers in order to simulate the calls from one caller component to the called one that has to be tested.\\
This continuous iteration of the bottom-up approach guarantees the testing coverage of the all the possible interactions of the components.\\
As previously mentioned, a thread analysis has to be performed too. This testing phase aims at verifying that the chains of function calls among components of different macro areas produce correct actions. The threads testing approach is chosen because it simulates the standard behaviour of the system, in terms of user requests. It could be considered a means to study the system performances too in this sense.\\
A final remark on the external components: the \textbf{MapService}, the \textbf{Payment System} and the \textbf{DBMS} components are already fully developed and in a bottom-up perspective they can be tested immediately using the corresponding system components as proper drivers. 
\section{Sequence of Component Integration}
%NOTE: The structure of this section may vary depending on the integration strategy you select in Section 2.3; use the structure proposed below as a non mandatory guide
This section will illustrate how the two different testing approaches, bottom-up and thread, are carried out. The order of presentation (that is bottom-up first and thread second) reflects the fact that to do a thread analysis one must have performed the bottom-up testing first.
A set of tables and figures will help to describe these two strategies.
\subsection{Bottom-Up Integration Strategy}
%For each subsystem, identify the sequence in which the software components will be integrated within the subsystem; relate this sequence to any product features that are being built up.
The bottom-up integration testing will be adopted within each \textit{Macro Area}. More precisely, for each pair of components that have function calls from one to the other a proper test suite will take care of testing all the possible interactions. After the testing of a pair of components, a third component which is in relation with one of this two tested modules is added and again its calls are checked. The same reasoning applies to the remaining components that belong to the same area.\\Let us introduce an example. Consider the components of the \textit{Input Area} as defined in the \textit{Elements to be Integrated} section: 
\imagebox{../img/ExampleBottomUp2}{Bottom-up test example}{0}{0.45}
\noindent where the following notation is adopted to state that Component 1 calls a function exposed by Component 2:
\imagebox{../img/ExampleBottomUp}{Bottom-up notation}{0}{0.45}
\noindent In the running example the bottom-up approach states that the already implemented part of the \textbf{ViewRender} \textit{drives} the already implemented portion of the \textbf{Dispatcher} component because the former calls all the methods exposed by the latter.\\The list of all the most significant methods invocations involved in the bottom-up testing for each pair of components of each \textit{Macro Area} is listed  in the \textit{Individual Steps and Test Description} section.
\subsection{Thread Integration Strategy}
%Identify the order in which subsystems will be integrated; if you have a single subsystem, 2.4.1 and 2.4.2 are to be merged in a single section.
The thread testing is performed in order to verify that chains of function calls among modules of different \textit{Macro Areas} lead to correct executions. This type of analysis is carried out when the components that it involves have a sufficient level of completion to support one specific functionality. The following figures show the most relevant threads analysis for the main functionalities of the system.\\Let us consider the \textit{Login Thread} as an example:
\imagebox{../img/Thread-LogIn}{Login Thread}{0}{0.8}
\noindent In order to complete the login operation the \textit{ViewRender} takes the request from the client and addresses it to the \textit{Dispatcher}. This last component realizes that this specific request should be handled by the \textit{LogInManager} which in turn needs to query the \textit{Model} to retrieve the client information.\\Please note that in these diagrams the purpose is only to show which components are needed to assess a specific functionality and not the entire control flow (see the sequence diagrams in the \textit{Design Document} for that). As such, some call links are omitted (for example the arc from the \textit{LogInManager} to the \textit{ViewRender} meaning that a web page is displayed at the end of the process is not present because it is not of interest).\\The following figures can be read with the same reasoning.
\imagebox{../img/Thread-RideStart}{Ride Start Thread}{0}{1}
\imagebox{../img/Thread-RideStop}{Ride Stop Thread}{0}{1}
\imagebox{../img/Thread-Reservation}{Reservation Thread}{0}{1}
\imagebox{../img/Thread-Registration}{Registration Thread}{0}{0.8}
\imagebox{../img/Thread-CarPlugIn}{Car Plug-in Thread}{0}{0.8}

\chapter[Individual Steps and Test Description]{Individual Steps and Test \\Description}
This section focuses on the interactions between pairs of components that will be progressively integrated. For each pair, a set of tests about the  function calls from one component to the other one is provided. This kind of test should cover all the possible calls in order to spot any type of undesirable behaviours just in time. For this reason, each function invocation is here evaluated many times under different circumstances depending on the actual values of the input parameters. Finally, for each such call the desired output is stated.
This integration test phase will be organized according to the logical areas division shown in the \textit{Elements to be Integrated} section.
For obvious space issues, in the current section only the most significant tests will be proposed, but keep in mind that such verification should be applied to every possible relation between the components.
\subsection{Management Area}
\subsubsection{Dispatcher$\,\to\,$ReservationManager}
\begin{tabularx}{\textwidth}{|X|X|}
	\hline
	\multicolumn{2}{|c|}{\textbf{ManageNewReservation(username,carID)}}\\
	\hline
	\textit{Input} & \textit{Result}\\
	\hline
	A null parameter & NullArgumentException\\
	\hline
	An empty or unknown parameter & InvalidArgumentException\\
	\hline
	Valid parameters & A new reservation for the given username is associated to the specified car\\
	\hline
\end{tabularx}
\subsubsection{Dispatcher$\,\to\,$RideManager}
\begin{tabularx}{\textwidth}{|X|X|}
	\hline
	\multicolumn{2}{|c|}{\textbf{StartRide(username,carID)}}\\
	\hline
	\textit{Input} & \textit{Result}\\
	\hline
	A null parameter & NullArgumentException\\
	\hline
	An empty or unknown parameter & InvalidArgumentException\\
	\hline
	Valid parameters & The RideManager registers that a new ride associated to the given username and car has started\\
	\hline
\end{tabularx}
\subsubsection{Dispatcher$\,\to\,$RideManager}
\begin{tabularx}{\textwidth}{|X|X|}
	\hline
	\multicolumn{2}{|c|}{\textbf{RideParams(saveMoneyOpt,finalDest)}}\\
	\hline
	\textit{Input} & \textit{Result}\\
	\hline
	A null parameter & NullArgumentException\\
	\hline
	An empty or unknown parameter & InvalidArgumentException\\
	\hline
	Valid parameters & The RideManager registers whether the user has enabled the Money Saving Option and his final destination\\
	\hline
\end{tabularx}
\subsubsection{Dispatcher$\,\to\,$RideManager}
\begin{tabularx}{\textwidth}{|X|X|}
	\hline
	\multicolumn{2}{|c|}{\textbf{RideStop(peopleOnBoard,position,batteryLevel)}}\\
	\hline
	\textit{Input} & \textit{Result}\\
	\hline
	A null parameter & NullArgumentException\\
	\hline
	An empty or unknown parameter & InvalidArgumentException\\
	\hline
	Valid parameters & The RideManager registers the number of peopleOnBoard, the final position and the remaining batteryLevel\\
	\hline
\end{tabularx}
\subsubsection{Dispatcher$\,\to\,$RideManager}
\begin{tabularx}{\textwidth}{|X|X|}
	\hline
	\multicolumn{2}{|c|}{\textbf{RidePayment(pluggedIn)}}\\
	\hline
	Valid parameter & The RideManager registers whether the user has plugged the car into the power grid\\
	\hline
\end{tabularx}
\subsubsection{Dispatcher$\,\to\,$StateManager}
\begin{tabularx}{\textwidth}{|X|X|}
	\hline
	\multicolumn{2}{|c|}{\textbf{ModifyCarState(carID,newState)}}\\
	\hline
	\textit{Input} & \textit{Result}\\
	\hline
	A null parameter & NullArgumentException\\
	\hline
	An empty or unknown parameter & InvalidArgumentException\\
	\hline
	Valid parameters & The StateManager updates the car with the given carID to the newState\\
	\hline
\end{tabularx}
\subsection{Input Area}
\subsubsection{ViewRender$\,\to\,$Dispatcher}
\begin{tabularx}{\textwidth}{|X|X|}
	\hline
	\multicolumn{2}{|c|}{\textbf{DispatchRequest(ReserveRequest)}}\\
	\hline
	\textit{Input} & \textit{Result}\\
	\hline
	A null parameter & NullArgumentException\\
	\hline
	An empty or unknown parameter & InvalidArgumentException\\
	\hline
	Valid parameter & The request is dispatched to the proper component\\
	\hline
\end{tabularx}
\subsubsection{ViewRender$\,\to\,$Dispatcher}
\begin{tabularx}{\textwidth}{|X|X|}
	\hline
	\multicolumn{2}{|c|}{\textbf{PickUpACar(username, carID)}}\\
	\hline
	\textit{Input} & \textit{Result}\\
	\hline
	A null parameter & NullArgumentException\\
	\hline
	An empty or unknown parameter & InvalidArgumentException\\
	\hline
	Valid parameters &  The ViewRender calls the suitable interface of the Dispatcher passing to it the input data\\
	\hline
\end{tabularx}
\subsubsection{Device$\,\to\,$ViewRender}
\begin{tabularx}{\textwidth}{|X|X|}
	\hline
	\multicolumn{2}{|c|}{\textbf{ReserveACar(username, carID)}}\\
	\hline
	\textit{Input} & \textit{Result}\\
	\hline
	A null parameter & NullArgumentException\\
	\hline
	An empty or unknown parameter & InvalidArgumetException\\
	\hline
	Valid parameters & The user inputs his username and the car he wants to reserve\\
	\hline
\end{tabularx}
\subsubsection{ViewRender$\,\to\,$Dispatcher}
\begin{tabularx}{\textwidth}{|X|X|}
	\hline
	\multicolumn{2}{|c|}{\textbf{ChangeCarState(carID, NewState)}}\\
	\hline
	\textit{Input} & \textit{Result}\\
	\hline
	A null parameter & NullArgumentException\\
	\hline
	An empty or unknown parameter & InvalidArgumentException\\
	\hline
	Valid parameters & The request of changing the state of the car with the specified carID is sent to the Dispatcher\\
	\hline
\end{tabularx}
\subsection{Ride Area}
\subsubsection{RideManager$\,\to\,$PaymentManager}
\begin{tabularx}{\textwidth}{|X|X|}
	\hline
	\multicolumn{2}{|c|}{\textbf{CheckBalance(username)}}\\
	\hline
	\textit{Input} & \textit{Result}\\
	\hline
	A null parameter & NullArgumentException\\
	\hline
	An empty or unknown parameter & InvalidArgumentException\\
	\hline
	Valid parameter & The user balance is returned\\
	\hline
\end{tabularx}
\subsubsection{RideManager$\,\to\,$MapController}
\begin{tabularx}{\textwidth}{|X|X|}
	\hline
	\multicolumn{2}{|c|}{\textbf{SearchSuggestedArea(FinalDestination)}}\\
	\hline
	\textit{Input} & \textit{Result}\\
	\hline
	A null parameter & NullArgumentException\\
	\hline
	An empty or unknown parameter & InvalidArgumentException\\
	\hline
	Valid parameter & The MapController computes the suggested areas where to park the car\\
	\hline
\end{tabularx}
\subsubsection{RideManager$\,\to\,$RideCostCalculator}
\begin{tabularx}{\textwidth}{|X|X|}
	\hline
	\multicolumn{2}{|c|}{\textbf{CalculateCost(peopleOnBoard,position,batterylevel)}}\\
	\hline
	\textit{Input} & \textit{Result}\\
	\hline
	A null parameter & NullArgumentException\\
	\hline
	An empty or unknown parameter & InvalidArgumentException\\
	\hline
	Valid parameters & The RideCostCalculator computes the total cost of the ride starting from the following input parameters: the number of people on board, the final position of the car, the final battery charge level\\
	\hline
\end{tabularx}
\subsubsection{RideManager$\,\to\,$MapController}
\begin{tabularx}{\textwidth}{|X|X|}
	\hline
	\multicolumn{2}{|c|}{\textbf{ChechPositon(MyPosition)}}\\
	\hline
	\textit{Input} & \textit{Result}\\
	\hline
	A null parameter & NullArgumentException\\
	\hline
	An empty or unknown parameter & InvalidArgumentException\\
	\hline
	Valid parameter &  The MapController checks the position of the car\\
	\hline
\end{tabularx}
\subsection{CarCommunication Area}
\subsubsection{ServerCommunicationManager$\,\to\,$CarCommunicationManager}
\begin{tabularx}{\textwidth}{|X|X|}
	\hline
	\multicolumn{2}{|c|}{\textbf{ReceiveReservation(ExpiringTime)}}\\
	\hline
	\textit{Input} & \textit{Result}\\
	\hline
	A null parameter & NullArgumentException\\
	\hline
	An empty or unknown parameter & InvalidArgumentException\\
	\hline
	Valid parameter &  The CarCommunicationManager receives the reservation and the expiring time for it\\
	\hline
\end{tabularx}
\subsubsection{CarCommunicationManager$\,\to\,$ServerCommunicationManager}
\begin{tabularx}{\textwidth}{|X|X|}
	\hline
	\multicolumn{2}{|c|}{\textbf{ReceiveRideStart(SaveMoneyOpt,FinalDest)}}\\
	\hline
	\textit{Input} & \textit{Result}\\
	\hline
	A null parameter & NullArgumentException\\
	\hline
	An empty or unknown parameter & InvalidArgumentException\\
	\hline
	Valid parameters &  The CarCommunicationManager notifies the ServerCommunicationManager that the ride is starting with indications about the preferences of the user\\
	\hline
\end{tabularx}
\subsubsection{ServerCommunicationManager$\,\to\,$CarCommunicationManager}
\begin{tabularx}{\textwidth}{|X|X|}
	\hline
	\multicolumn{2}{|c|}{\textbf{CommunicateParkArea(AreaPosition)}}\\
	\hline
	\textit{Input} & \textit{Result}\\
	\hline
	A null parameter & NullArgumentException\\
	\hline
	An empty or unknown parameter & InvalidArgumentException\\
	\hline
	Valid parameter & The ServerCommunicationManager communicates to the CarCommunicationManager the predefined position of the area where the user can park the car\\
	\hline
\end{tabularx}
\subsubsection{CarCommunicationManager$\,\to\,$ServerCommunicationManager}
\begin{tabularx}{\textwidth}{|X|X|}
	\hline
	\multicolumn{2}{|c|}{\textbf{RideStop(carID,peopleOnBoard,position,batteryLevel)}}\\
	\hline
	\textit{Input} & \textit{Result}\\
	\hline
	A null parameter & NullArgumentException\\
	\hline
	An empty or unknown parameter & InvalidArgumentException\\
	\hline
	Valid parameters & The CarCommunicationManager notifies that the user has ended the ride on the specified car. Information about the peopleOnBoard, the final position and the batteryLevel are also provided\\
	\hline
\end{tabularx}
\subsubsection{ServerCommunicationManager$\,\to\,$CarCommunicationManager}
\begin{tabularx}{\textwidth}{|X|X|}
	\hline
	\multicolumn{2}{|c|}{\textbf{SendCost(Cost)}}\\
	\hline
	\textit{Input} & \textit{Result}\\
	\hline
	A null parameter & NullArgumentException\\
	\hline
	An empty or unknown parameter & InvalidArgumentException\\
	\hline
	Valid parameter & The ServerCommunicationManager sends the cost of the ride to the CarCommunicationManager\\
	\hline
\end{tabularx}
\subsubsection{CarCommunicationManager$\,\to\,$ServerCommunicationManager}
\begin{tabularx}{\textwidth}{|X|X|}
	\hline
	\multicolumn{2}{|c|}{\textbf{SendPlugInTimeout(PluggedIn)}}\\
	\hline
	\textit{Input} & \textit{Result}\\
	\hline
	A null parameter & NullArgumentException\\
	\hline
	An empty or unknown parameter & InvalidArgumentException\\
	\hline
	Valid parameter & The CarCommunicationManager indicates if the user has plugged the car into the power grid\\
	\hline
\end{tabularx}
\subsection{Car Area}
\subsubsection{CentralUnit$\,\to\,$CarCommunicationManager}
\begin{tabularx}{\textwidth}{|X|X|}
	\hline
	\multicolumn{2}{|c|}{\textbf{SendBackRideStart(SaveMoneyOpt,FinalDestination)}}\\
	\hline
	\textit{Input} & \textit{Result}\\
	\hline
	A null parameter & NullArgumentException\\
	\hline
	An empty or unknown parameter & InvalidArgumentException\\
	\hline
	Valid parameters & The CentralUnit sends to the CarCommunicationManager the preferences of the user in terms of the money saving option and the final destination of the ride.\\
	\hline
\end{tabularx}
\subsubsection{CarCommunicationManager$\,\to\,$CentralUnit}
\begin{tabularx}{\textwidth}{|X|X|}
	\hline
	\multicolumn{2}{|c|}{\textbf{StoreParkPosition(AreaPosition)}}\\
	\hline
	\textit{Input} & \textit{Result}\\
	\hline
	A null parameter & NullArgumentException\\
	\hline
	An empty or unknown parameter & InvalidArgumentException\\
	\hline
	Valid parameter & The CarCommunicationManager sends to the CentralUnit the position of the area where the user can park\\
	\hline
\end{tabularx}
\subsubsection{CentralUnit$\,\to\,$ScreenManager}
\begin{tabularx}{\textwidth}{|X|X|}
	\hline
	\multicolumn{2}{|c|}{\textbf{DisplayParkPosition(Position)}}\\
	\hline
	\textit{Input} & \textit{Result}\\
	\hline
	A null parameter & NullArgumentException\\
	\hline
	An empty or unknown parameter & InvalidArgumentException\\
	\hline
	Valid parameter & The ScreenManager displays on the screen the position on the map where the user can park to obtain special discount\\
	\hline
\end{tabularx}
\subsubsection{CentralUnit$\,\to\,$CarCommunicationManager}
\begin{tabularx}{\textwidth}{|X|X|}
	\hline
	\multicolumn{2}{|c|}{\textbf{HandleStop(CarID,PeopleOnBoard,Position,BatteryLevel)}}\\
	\hline
	\textit{Input} & \textit{Result}\\
	\hline
	A null parameter & NullArgumentException\\
	\hline
	An empty or unknown parameter & InvalidArgumentException\\
	\hline
	Valid parameters & The CarCommunicationManager receives from the CentralUnit all the data that have to be passed the to system in order to properly manage the end of the ride\\
	\hline
\end{tabularx}
\subsubsection{CentralUnit$\,\to\,$CarCommunicationManager}
\begin{tabularx}{\textwidth}{|X|X|}
	\hline
	\multicolumn{2}{|c|}{\textbf{PlugInTimeout(PluggedIn)}}\\
	\hline
	\textit{Input} & \textit{Result}\\
	\hline
	A null parameter & NullArgumentException\\
	\hline
	An empty or unknown parameter & InvalidArgumentException\\
	\hline
	Valid parameter & The CentralUnit notifies the CarCommunicationManager whether the user has plugged the car into the power grid in time\\
	\hline
\end{tabularx}
\subsection{Render Area}
\subsubsection{MapController$\,\to\,$ViewRender}
\begin{tabularx}{\textwidth}{|X|X|}
	\hline
	\multicolumn{2}{|c|}{\textbf{ShowAvailableCars(position,range)}}\\
	\hline
	\textit{Input} & \textit{Result}\\
	\hline
	A null parameter & NullArgumentException\\
	\hline
	An empty or unknown parameter & InvalidArgumentException\\
	\hline
	Valid parameters & A webpage with the available cars within the range of distance from the position specified is displayed by the ViewRender\\
	\hline
\end{tabularx}
\subsubsection{LogInController$\,\to\,$ViewRender}
\begin{tabularx}{\textwidth}{|X|X|}
	\hline
	\multicolumn{2}{|c|}{\textbf{ShowMainPage(username)}}\\
	\hline
	\textit{Input} & \textit{Result}\\
	\hline
	A null parameter & NullArgumentException\\
	\hline
	An empty or unknown parameter & InvalidArgumentException\\
	\hline
	Valid parameter & The Main Page of the specified user is displayed by the ViewRender\\
	\hline
\end{tabularx}
\subsubsection{RideManager$\,\to\,$ViewRender}
\begin{tabularx}{\textwidth}{|X|X|}
	\hline
	\multicolumn{2}{|c|}{\textbf{AbortPickUp(errorMsg)}}\\
	\hline
	\textit{Input} & \textit{Result}\\
	\hline
	A null parameter & NullArgumentException\\
	\hline
	An empty or unknown parameter & InvalidArgumentException\\
	\hline
	Valid parameter & A webpage with the critical error is displayed by the ViewRender\\
	\hline
\end{tabularx}
\subsection{Data Area}
\subsubsection{RideManager$\,\to\,$Model}
\begin{tabularx}{\textwidth}{|X|X|}
	\hline
	\multicolumn{2}{|c|}{\textbf{ChangeCarState(carID,newState)}}\\
	\hline
	\textit{Input} & \textit{Result}\\
	\hline
	A null parameter & NullArgumentException\\
	\hline
	An empty or unknown parameter & InvalidArgumentException\\
	\hline
	Valid parameter & The state of the car with carID is set to newState on the database\\
	\hline
\end{tabularx}
\subsubsection{RegistrationManager$\,\to\,$Model}
\begin{tabularx}{\textwidth}{|X|X|}
	\hline
	\multicolumn{2}{|c|}{\textbf{InsertNewUser(credentials,username,licenseNumber,email,paymentinfo}}\\
	\hline
	\textit{Input} & \textit{Result}\\
	\hline
	A null parameter & NullArgumentException\\
	\hline
	An empty or unknown parameter & InvalidArgumentException\\
	\hline
	Valid parameters & A new record for a new user is created in the Model\\
	\hline
\end{tabularx}
\subsubsection{LogInManager$\,\to\,$Model}
\begin{tabularx}{\textwidth}{|X|X|}
	\hline
	\multicolumn{2}{|c|}{\textbf{FindUser(username,password)}}\\
	\hline
	\textit{Input} & \textit{Result}\\
	\hline
	A null parameter & NullArgumentException\\
	\hline
	An empty or unknown parameter & InvalidArgumentException\\
	\hline
	Valid parameters & The LogInManager checks that user is already registered\\
	\hline
\end{tabularx}
%For each step of the integration process above, describe the type of tests that will be used to verify that the elements integrated in this step perform as expected. Describe in general the expected results of the test set. You may refer to Chapter 3 and Chapter 4 of the test plan example [1] as an example of what we expect.
%(NOTE: This is not a detailed description of test protocols. Think of this as the test design phase. Specific protocols will be written to fulfill the goals of the tests in this section.)
\chapter{Tools and Test Equipment Required}
%Identify all tools and test equipment needed to accomplish the integration. Refer to the tools presented during the lectures. Explain why and how you are going to use them.
%Note that you may also use manual testing for some part. Consider manual testing as one of the possible tools you have available.
\section{Tools}
Testing tools are important because they can automatize and speed up the testing of the system; the used tools that will be presented later in this section, are then useful to properly test firstly the single components and secondly how they integrate and communicate with each other in order to form the complete system and perform all the operations.\\
Considering that the chosen technology for the implementation of webapplication is JavaEE, the following testing tools are taken into account:
\begin{itemize}
	\item Mockito and JUnit are used in the first step of this process because, nevertheless they are mainly designed for unit testing, using them it is possible to exploit all the dependencies and interactions for each couple of components, as done in the bottom-up phase integration for the component in each macroarea but also while integrating the full 'chain' of components that realize a single thread
	\item Arquillian framework because it lets write test cases against Java Containers and also to check dependency injection of components and transaction control (for example for the Data macroarea that involves interactions with the DBMS).
\end{itemize}

\section{Test Equipment}
In order to properly execute the whole integration testing phase on the system, different type of equipments are required to carry out both the testing on the purely software part (that will be deployed on the server) and the part which will be embedded into the car (that includes the on-board touch screen, the central unit and the specific component devoted to the communication with the server). \\
In particular, starting then from the previous statement, it is fundamental to have:
\begin{itemize}
	\item at least a device for each type (personal computer, tablet and smart-phone) that will be used from the user to access the system's functionalities and from the employees of the car-sharing company
	\item the application server, properly configured and already working with the webapplication deployed on it
	\item the DBMS already initialized (with all the tables required) and running on a specific server machine
	\item both the external component (that is the MapService and the PaymentSystem) must be available through their APIs
	\item a fully working prototype of the subsystem that will be installed in each car, complete with the communication part.
	\item a fully working prototype of a charging station that will be available in each charging area
\end{itemize}
For a better and real testing context it would better to have a complete real car and charging station too. However, while testing for the first time the system it is fine to have only a simplified model of both car and charging station because it can be costly and quite difficult to have nearby the development area but when the testing is in an advanced phase real cars and charging stations should be used to have a better feedback from the system.\\
In conclusion, for constructing a simply environment to properly test the overall system, all these elements are the minimum that is required; after having checked the all the functionalities are working properly for this simple case, it is possible to add more components (for example more devices, cars and/or charging stations) to simulate a more realistic case that let's also to check the performance and the reliability of the developed system.

\chapter[Program Stubs and Test Data Required]{Program Stubs and Test Data\\Required}
Using a bottom-up approach to test each pair of every macro-area previously defined, and a thread strategy to test the multiple iterations of components on a single action, we have to define some drivers that can simulate the invocations even if the components are fully developed.
Here is the list of the drivers that will be developed in the integration testing phase.
\begin{itemize}
	\item \textbf{ReservationManager Driver}: is used to simulate methods that used in reservation phase
	\item \textbf{RegistrationManager Driver}: is built in order to expose methods which you can use to register with.
	\item \textbf{StateManager Driver}: provides methods that are used to change the state of the car
	\item \textbf{LogInManager Driver}: exposes interfaces to test the login into the system
	\item \textbf{RideManager Driver}: is used to give methods that are used to manage the ride
\end{itemize}
These drivers are used to link the \textit{Dispatcher} to the \textit{Model}, in order to simulate the communication between each component of the \textit{Management Area}
\begin{itemize}
	\item \textbf{RideCostCalculator Driver}: is built to provide methods used for the calculation of the cost of the ride
\end{itemize}
This is used in order to simulate a true communication between the \textit{Ride Area} and the \textit{Input Area}
\begin{itemize}
	\item \textbf{ServerCommunicationManager Driver}: is a server component used to simulate the most important iteration between the server to the car
	\item \textbf{CarCommunicationManager Driver}: is a car interface used to provide methods that are used to give a communication between car and server
\end{itemize}
These are used to link the car communication part to the server one, creating a connection between the \textit{CarCommunication Area} and the \textit{Input Area}.
\begin{itemize}
	\item \textbf{Model Driver}:is used to invoke methods that are linked to the \textbf{DBMS}
\end{itemize}
This is used to link the \textit{Input Area} to the external \textit{DBMS}.
As you can see, there are no need to build a \textbf{Dispatcher Driver}, or a \textbf{PaymentManager Driver} or a \textbf{MapController Driver} because, as it was said previously, they are fully developed because they have to manage request or interact to external components.
%Based on the testing strategy and test design, identify any program stubs or special test data required for each integration step.
\chapter{Effort Spent}
In order to complete this document, each author worked for 15 hours.
%-------------------------------------------------------------------------
%   END DOCUMENT
%-------------------------------------------------------------------------
\end{document} 